\sscp{Collins}{07-01-02}

In many publications Collins has investigated the subgrouping
of the languages of geographical central Maluku.
An attempt is made in \frf{07-07} to combine
Collins’ (\citeyear{Collins1981}, \citeyear[20--21, 102]{Collins1983},
\citeyear[127]{Collins1986}) higher level classifications of these languages
in order to provide context for the following
discussion.\footnote{
		\citet[20]{Collins1983} places Seti as a
		fifth branch under his \tsc{Nunusaku},
		but in \citet[127]{Collins1986} Seti is
		reclassified to be a branch under East Seram.
		In the \citeyear{Collins1983} work,
		Proto-East Central Maluku has two branches; in the \citeyear{Collins1986} work
		it is revised to have three branches.
		These later revisions are reflected in \frf{07-07}.}

\begin{figure}[h]\centering
	\includegraphics[width=0.685\columnwidth]{CollinsTree}
	\fcp{Summary of Collins’ (1981, 1983, 1986) classification of central Maluku languages}{07-07}
\end{figure}

%\dirtree{%The percent after { is needed
%.1 {Proto-Central Maluku (PCM)}.
		%.2 {Proto-West Central Maluku}.
				%.3 {Buru-Sula-Taliabo}.
						%.4 {Sula-Taliabo}.
						%.4 {Buru}.
				%.3 {Ambelau}.
		%.2 {Proto-East Central Maluku}.
				%.3 {Nunusaku}.
						%.4 {Kayeli}.
						%.4 {Three Rivers}.
						%.4 {Piru Bay}.
						%.4 {Patakai-Manusela}.
				%.3 {East Seram}.
						%.4 {Seti}.
						%.4 {East Rivers}.
				%.3 {Proto-Banda}.
						%.4 {Banda}.
						%.4 {Seran Laut}.
%}

\citeauthor{Collins1986}’ (\citeyear[134]{Collins1986}) Proto-Banda, which includes Banda
and our \tsc{Eastern Islands} languages of \tsc{Seram-Tanimbar-Bomberai},
is based on ten retentions,
and one innovation *d/*z/*R > **r.
The merger of *d/*z/*R is found more widely in several branches
of our \tsc{Seram-Tanimbar-Bomberai}, and is problematic
for \tsc{Eastern Islands} when data from additional
languages is considered (\srf{11-03}, \srf{12-02}).
All up, the collective differences in the historical phonology
and the lexicon and morphosyntax between Banda
and \tsc{Eastern Islands} lead us to conclude that the similarity
of PMP *R > \it{r} in Geser-Gorom
and Banda is better explained as parallel development,
particularly given that it may be a retention depending on the
phonetic interpretation of *R.
This is discussed in more detail in \srf{12-02}.

At a higher level,
what is used by Collins to link West Central Maluku with Nunusaku
does not overlap with what is used to link East Seram
and Proto-Banda languages with Nunusaku,
so there are problems trying to link the two peripheral groups
together under a single Central Maluku subgroup.
\citet[40]{Collins1981} lists the innovations that link his West
Central Maluku (our \tsc{Taliabu}
and \tsc{Sula-Buru}) with his East Central Maluku (our \tsc{Ambon-Seram},
two nodes of \tsc{Seram-Tanimbar-Bomberai},
and Banda) under his Central Maluku.
However, updating to modern understanding of MP historical phonology,
we are left with five retentions,
two non-exclusive mergers, and one error,
none of which meet the criteria for defining subgroups.\footnote{
		Modern notation is in parentheses.
		Five would now be considered retentions,
		not innovations: *z/*Z > *Z (*z = **z); *d/*D > *ḍ (*d = **d);
		*-ay/*-ey > *-ay (*-ay = **-ay); *S > *S (\tsc{PAn} *S = *S,
		PMP *h = **h); *q > *q (*q = **q).
		Two are not exclusive: *mp/*mb > *mb; *nd/*nt > *nd (consistent
		with \citeauthor{Blust1981}'s (\citeyear[26]{Blust1981})
		claims for the wider “CMP” region).
		One is an error: *j > \it{t}/{\gap}{\#},
		since this final \it{-t} is actually a productive nominal suffix
		in languages across the region,
		and not a reflex of *j.
		In \citet{Collins1981,Collins1983} the claim for *j > \it{t}{\gap}{\#} was
		focused on one etymon: *waRəj ‘vine,
		rope’, whose reflexes also occur without \it{t}{\#},
		and also with other suffixes in several languages.
		(see footnote \ref{fn:Collins-Final-t}, this chapter.)}
Our investigation finds no
exclusively shared innovations in the historical phonology that
support Collins’ higher level Central Maluku subgroup.
They do share the non-exclusive innovation of *-ay > \it{a}.

\citet{Collins1983} provides an internal classification of the
languages of the Ambon-Lease islands,
along with those of western
and central Seram and the Manipa Strait between Seram and Buru.
The \citeyear{Collins1983} work focuses primarily on only two branches of his Nunusaku
subgroup (\tsc{Three Rivers} [= our \tsc{Atamanu} and \tsc{Wemale}] and \tsc{Piru Bay}),
and does not systematically discuss what innovations might define Nunusaku.
Because his focus is elsewhere,
\citet{Collins1983} does not define criteria
for the \tsc{Patakai-Manusela} branch of Nunusaku.
Collins’ Patakai-Manusela is very roughly equivalent to Stresemann’s
undefined Sub-Seran.